\chapter{Rastreabilidade da versão documental}
\label{ap:rastreabilidade}
\label{ap:reproducao}
\section{Origem dos resultados}
As tabelas e os gráficos desta monografia derivam da execução identificada como \texttt{fechamento\_20260928\_cuda}. O código científico correspondente foi registrado no commit \texttt{91a2a28360ded41fd86abf3371c989d571bf95bd}. A redação iniciou sobre o commit documental \texttt{9894a6a325ab3b155cbe52fa4c6aa78fd47e2db2}. Esses identificadores não significam que esta versão textual tenha sido aprovada ou publicada.

O manifesto final original teve seu resumo SHA-256 conferido (sequência contínua, apresentada em duas linhas):
\begin{center}\small\ttfamily
5fc6acd8c42c75520c0eecefa5366f3a\\
33e0627d217b0b0cf61dd556876f34c9
\end{center} Um inventário local registra os hashes de 72 arquivos, totalizando 81.607.118 bytes. Esse inventário permite detectar alterações no conjunto preservado, mas não representa um backup independente.

\begin{table}[htbp]
\centering\caption{Relação entre artefatos e elementos da monografia.}
\begin{tabularx}{\textwidth}{>{\raggedright\arraybackslash}p{.34\textwidth}>{\raggedright\arraybackslash}X}\toprule
Artefato de origem & Uso documental\\\midrule
\texttt{fusion\_metrics.csv} & Tabela principal e gráfico de mAP no Gallagher.\\
\texttt{fusion\_metrics\_}\allowbreak\texttt{per\_query.csv} & Conferência das médias, diferenças pareadas e contagens de ganhos, empates e perdas.\\
\texttt{face\_metrics.csv} & Resultado auxiliar de recuperação em LFW.\\
\texttt{global\_metrics.csv} & Resultado auxiliar em Holidays.\\
Manifestos e protocolo & Populações, seleção de consultas, modelos, dispositivos e regras de avaliação.\\
Fechamento técnico & Testes e aceitação operacional previamente registrados.\\\bottomrule
\end{tabularx}
\fonte{organização documental do autor.}
\end{table}

\section{Pacote leve e alcance da conferência}
O diretório documental de evidências contém cópias dos três arquivos agregados, um resumo pareado sem identificadores de consulta e um registro de procedência derivado. O manifesto derivado não substitui nem se apresenta como o manifesto original. O arquivo \texttt{PROVENIENCIA\_FIGURAS\_TABELAS.json} registra hashes de entrada e saída dos elementos gerados. A rotina documental confere que as 100 linhas correspondem a cinco configurações e vinte consultas, com 588 candidatos e 884 relações de relevância em cada configuração.

Não são incluídos no pacote de divulgação fotografias, recortes, representações biométricas, pesos de modelos ou rankings que permitam vincular pessoas às consultas. Essa seleção de conteúdo não é uma alegação formal de anonimização. Os arquivos detalhados continuam sujeitos às condições de acesso e de uso da pesquisa. Como apenas os dez primeiros resultados foram persistidos, as médias de AP podem ser conferidas contra os valores por consulta, mas o cálculo integral de AP não pode ser refeito apenas com os rankings salvos.

\section{Correspondência com a proposta inicial}
A comparação abaixo resume os eixos da proposta assinada consultada, sem alterar seu texto ou pressupor aprovação institucional das mudanças. O detalhamento documental permanece em \texttt{MAPA\_EVIDENCIAS.md}.

\begin{table}[htbp]
\centering\small\caption{Realização e limites dos objetivos da proposta.}
\begin{tabularx}{\textwidth}{p{.31\textwidth}X}\toprule
Eixo proposto & Correspondência no trabalho realizado\\\midrule
a) Organizar bases e acervo & Bases existentes e anotações originais; família somente para aceitação privada.\\
b) Comparar e implementar algoritmos faciais & SCRFD/ArcFace integrados; sem comparação entre detectores ou reconhecedores.\\
c) Descritores globais com CNN & ResNet50 previamente treinada, sem treinamento novo.\\
d) FAISS e grande escala & Busca vetorizada existente; FAISS e ensaio de escala não realizados.\\
e) Fusão e interface amigável & Pesos fixos e interface com gestão de álbuns; sem estudo de usabilidade.\\
f) Métricas e latência & P@K, R@K e mAP; sem estudo controlado de latência.\\\bottomrule
\end{tabularx}
\fonte{síntese da proposta consultada e dos artefatos atuais; eixos parafraseados.}
\end{table}

\section{Revisão autoral e uso de assistência}
A assistência de inteligência artificial está declarada na nota inicial. Não houve nova inferência nem geração de resultados sintéticos. A compreensão, revisão e aprovação do texto pelo autor permanecem pendentes, assim como a confirmação da forma institucional de declaração desse uso.
